\chapter*{Словарь терминов}             % Заголовок
\addcontentsline{toc}{chapter}{Словарь терминов}  % Добавляем его в оглавление

\textbf{Информационно-поисковая система (ИПС)}~--- программно-аппаратная
система, обеспечивающая обнаружение, сопоставление, идентификацию и~поиск
объектов в~массиве визуальных данных (изображений и~видеопотоков) при
заданных эксплуатационных ограничениях по~времени отклика,
вычислительным ресурсам и~метрикам качества.

\textbf{Многокомпонентная структура}~--- объект, описываемый совокупностью
топологически и~семантически связанных компонентов $C_1,\ldots,C_K$,
наблюдаемых независимо и~допускающих самостоятельное признаковое
описание. В~частном случае~$K=2$ многокомпонентной структурой является
человек, описываемый связанной парой компонентов <<силуэт-лицо>>;
в~смежных предметных областях~--- транспортное средство (<<кузов~---
регистрационный знак>>), металлоконструкция (<<участок~--- локальный
дефект>>), беспилотный летательный аппарат и~группа малоразмерных
объектов и~др.

\textbf{Детектирование}~--- процесс автоматического выделения областей
изображения или видеопотока, соответствующих объектам заданного класса
или компонентам многокомпонентной структуры; формализуется как функция
$D(\cdot)$, возвращающая множество ограничивающих рамок (bounding boxes)
для входного изображения.

\textbf{Связанные объекты <<силуэт-лицо>>}~--- частный случай связанной
пары компонентов многокомпонентной структуры (при~$K=2$), в~которой
лицо локализовано внутри ограничивающей рамки силуэта; согласованное
детектирование таких пар~--- ключевое требование к~первому этапу
методики настоящей работы.

\textbf{Реидентификация (person re-identification, ReID)}~--- задача
поиска и~распознавания одного и~того же субъекта на~различных
изображениях или видеозаписях, полученных с~непересекающихся камер
видеонаблюдения, в~условиях изменения освещения, ракурса, позы,
частичных окклюзий и~вариативности внешнего вида.

\textbf{Эмбеддинг (embedding, признаковое представление)}~--- вектор
$\mathbf{z}_k\in\mathbb{R}^{d_k}$ фиксированной размерности, формируемый
кодировщиком~$\Phi_k$ как компактное и~дискриминативное численное
описание $k$-го компонента наблюдаемого объекта.

\textbf{ROI (region of interest, область интереса)}~--- ограниченная
прямоугольная область на~изображении, соответствующая выделенному
детектором компоненту структуры; служит входом для кодировщика
эмбеддингов.

\textbf{Кодировщик (encoder)}~--- нейросетевая модель, отображающая
ROI компонента в~эмбеддинг; в~настоящей работе кодировщики реализованы
на~базе архитектуры Vision Transformer (ViT), обучаемой с~функцией потерь
ArcFace.

\textbf{Vision Transformer (ViT)}~--- архитектура нейросетевого
кодировщика изображений, основанная на~разбиении входа на~непересекающиеся
патчи и~их обработке стеком блоков механизма самовнимания (self-attention).

\textbf{ArcFace}~--- функция потерь метрического обучения, вводящая
аддитивную угловую margin-поправку в~классификационную границу
и~обеспечивающая формирование компактных, разделимых межклассовых
кластеров эмбеддингов в~признаковом пространстве.

\textbf{YOLOv8}~--- современная нейросетевая архитектура объектного
детектирования на~базе блоков C2f и~модуля Spatial Pyramid Pooling Fast
(SPPF), обеспечивающая высокую точность локализации при работе в~режиме,
близком к~реальному времени.

\textbf{Каскадный эффект}~--- накопление и~распространение ошибок
последовательных этапов обработки видеопотока (детектирование~$\to$
извлечение эмбеддингов~$\to$ интеграция~$\to$ принятие решения),
при котором неточности раннего этапа задают верхнюю границу качества
последующих.

\textbf{Информационная неопределённость}~--- эксплуатационные условия
ИПС, характеризующиеся временной недоступностью отдельных компонентов
структуры (например, отсутствием фронтального лица), вариативностью
внешнего вида объекта, окклюзиями, шумами и~различиями условий съёмки
между камерами.

\textbf{Байесовская интеграция}~--- метод объединения признаковых
описаний компонентов структуры на~основе теоремы Байеса:
апостериорная вероятность класса вычисляется как произведение
правдоподобий по~доступным компонентам и~априорной вероятности.

\textbf{MAP-критерий (Maximum A~Posteriori)}~--- правило принятия
решения по~максимуму апостериорной вероятности:
$y^{*}=\arg\max_{y\in Y} P(y\mid\{\mathbf{z}_k\colon k\in\mathcal{M}(t)\})$.
Эквивалентно максимизации суммы логарифмов правдоподобий по~доступным
компонентам.

\textbf{Гауссова модель правдоподобия}~--- аппроксимация распределения
эмбеддингов класса в~признаковом пространстве компонента многомерным
нормальным распределением $\mathcal{N}(\boldsymbol{\mu}_{y,k},\Sigma_{y,k})$.

\textbf{Ковариационная матрица класса}~--- параметр гауссовой модели
$\Sigma_{y,k}$, отражающий разброс эмбеддингов класса~$y$ внутри
признакового пространства компонента~$k$; в~предлагаемой методике
играет роль коэффициента доверия модальности при автоматическом
взвешивании компонентов в~MAP-критерии.

\textbf{Расстояние Махаланобиса}~--- метрика
$d_{\mathrm{M}}(\mathbf{z};\boldsymbol{\mu},\Sigma)=
\sqrt{(\mathbf{z}-\boldsymbol{\mu})^{\mathsf{T}}\Sigma^{-1}(\mathbf{z}-\boldsymbol{\mu})}$,
учитывающая ковариационную структуру распределения и~входящая
в~логарифм гауссова правдоподобия.

\textbf{$L_2$-нормировка}~--- приведение вектора к~единичной евклидовой
длине: $\hat{\mathbf{z}}=\mathbf{z}/\lVert\mathbf{z}\rVert_2$.
Применяется к~эмбеддингам после извлечения для устойчивости последующих
сравнений и~для совместимости с~ArcFace и~байесовской интеграцией.

\textbf{Атрибутивные признаки}~--- семантические характеристики объекта
(одежда, аксессуары, цветовая палитра, ракурс, контекст сцены),
дополняющие лицевые и~силуэтные эмбеддинги в~многомодальных схемах
реидентификации.

\textbf{Score-level fusion (объединение на~уровне оценок)}~---
стратегия мультимодального объединения, при~которой решения, принятые
по~отдельным модальностям, агрегируются на~уровне числовых оценок
сходства; противопоставляется feature-level fusion (объединению
на~уровне признаков).

\textbf{Re-ranking (переранжирование)}~--- этап постобработки в~задаче
реидентификации, в~котором результаты первичного поиска корректируются
с~учётом структуры ближайших соседей в~признаковом пространстве и/или
взаимных сходств кандидатов.

\textbf{Метрики качества реидентификации}: \emph{Precision}~---
точность, доля правильных положительных решений; \emph{Recall}~---
полнота, доля найденных истинных совпадений; $F_1$-мера~--- гармоническое
среднее точности и~полноты; \emph{mAP} (mean Average Precision)~---
среднее значение средней точности по~запросам;
\emph{CMC/Rank-$k$}~--- кумулятивная характеристика правильного
попадания в~top-$k$ кандидатов; \emph{ROC/AUC}~---
характеристическая кривая и~площадь под~ней.
